\section{Path diagrams and missing recurrent classes}
\label{sec:diagrams}

We first describe finite paths through a segment, without assuming that
the paths come from an automaton. Our goal is a loss measure that detects
edge inclusion and remains controlled when a segment is put in context.

Relations are composed in \emph{path order}: $x(AB)y$ means that
$xAz$ and $zBy$ for some $z$. On a single set, $A^*$ denotes the reflexive
transitive closure, so it includes paths of length zero. We use the
elementary path identities
\begin{equation}\label{diag:star}
 (A\cup B)^*=A^*(BA^*)^*,\qquad
 A(BA)^*=(AB)^*A.
\end{equation}
The first groups successive $A$-steps between $B$-steps; the second groups
the same alternating path from its opposite end. These identities hold
whenever the relation types permit the displayed products.

\subsection{The diagram monoid}

Fix disjoint sets $D^+$ and $D^-$, each of size $m\ge1$.
A diagram $a$ consists of four arbitrary relations
\[
 F_a\subseteq D^+\times D^+,\quad
 B_a\subseteq D^-\times D^-,\quad
 L_a\subseteq D^+\times D^-,\quad
 R_a\subseteq D^-\times D^+.
\]
Entering the segment on the left uses a label of $D^+$; $F_a$ records an
exit on the right and $L_a$ an exit on the left. Entering on the right uses
a label of $D^-$; $B_a$ records an exit on the left and $R_a$ an exit on
the right. The relations $F_a,B_a$ are the \emph{through relations}, and
$L_a,R_a$ are the \emph{return relations}; see Figure~\ref{fig:diagram}.

\begin{figure}[tb]
\centering
\begin{tikzpicture}[>=Stealth,thick,font=\small]
 \fill[black!3] (0.4,-0.35) rectangle (6.6,2.35);
 \draw[dashed,black!45] (0.4,-0.4)--(0.4,2.4);
 \draw[dashed,black!45] (6.6,-0.4)--(6.6,2.4);
 \node (lp) at (1,2) {$D^+$};
 \node (lm) at (1,0) {$D^-$};
 \node (rp) at (6,2) {$D^+$};
 \node (rm) at (6,0) {$D^-$};
 \draw[->] (lp)--node[above] {$F_a$}(rp);
 \draw[->] (rm)--node[below] {$B_a$}(lm);
 \draw[->] (lp) to[bend right=35] node[left] {$L_a$}(lm);
 \draw[->] (rm) to[bend right=35] node[right] {$R_a$}(rp);
 \node at (3.5,1) {segment $a$};
 \node[anchor=east] at (-0.15,2) {entry};
 \node[anchor=east] at (-0.15,0) {exit};
 \node[anchor=west] at (7.15,2) {exit};
 \node[anchor=west] at (7.15,0) {entry};
\end{tikzpicture}
\caption{The four relations of a path diagram. Each arrow represents a
relation between entire label sets, not a single deterministic edge.
In a product, an exit enters the neighboring segment with the same label.}
\label{fig:diagram}
\end{figure}

Place diagrams in a line. At each internal boundary, identify an exit
with the entry of the neighboring diagram bearing the same label.
The product records finite paths from an outer entry to the first outer
exit, with all earlier joins internal. A traversal of any consecutive
block can be replaced by one edge of the block's product, and that edge
can conversely be expanded into a finite traversal. This proves
associativity. The diagram with $F$ and $B$ identity relations and $L,R$
empty is the identity. Thus the diagrams form a finite monoid
$\mathcal T_m$.

This is the degree-$m$ monoid of partitioned binary relations of
Martin and Mazorchuk \cite[Sections~2.1--2.3]{MartinMazorchuk2013}.
To identify the two descriptions, label $D^+$ and $D^-$ by a common
$m$-element set. At each boundary, an entry and an exit with the same
underlying label give one vertex. The four relations above become
the four blocks of a binary relation on the left and right vertex
sets. A join always switches from one factor to its neighbor, exactly
as in their alternating-path composition. We have included the
path expansion and contraction proof of associativity to fix the
finite-path convention used below.

Write $a\sqsubseteq b$ when all four relations of $a$ are included in
the corresponding relations of $b$. Products preserve this order in
each argument. Joining two diagrams at one boundary gives
\begin{equation}\label{diag:product}
\begin{aligned}
 F_{ab}&=F_a(L_bR_a)^*F_b,&
 B_{ab}&=B_b(R_aL_b)^*B_a,\\
 L_{ab}&=L_a\cup F_a(L_bR_a)^*L_bB_a,&
 R_{ab}&=R_b\cup B_bR_a(L_bR_a)^*F_b.
\end{aligned}
\end{equation}
For example, a forward path crosses $a$, alternates returns at the join,
and crosses $b$. Reflection exchanges $D^+$ with $D^-$, $F$ with $B$,
and $L$ with $R$, while reversing product order. We will prove symmetric
statements for the positive sign and obtain the negative sign by this
reflection.

\subsection{From complementation to an order-reversing image}

The algebraic obstruction to complementation is the following bound.
As in the introduction, $\mathcal R_H$ is the monoid of all binary
relations on $H$, composed in path order.

\begin{theorem}[Order-reversing image bound]
\label{amp:image}
\label{amp:relation-bound}
Let $m\geq1$, let $H$ be a finite set of cardinality $h\geq2$, and let
$S_0$ be a submonoid of $\mathcal T_m$.  Suppose that a surjective
multiplicative map $\phi_0:S_0\to\mathcal R_H$ satisfies
\begin{equation}
 z\sqsubseteq w\quad\Longrightarrow\quad
 \phi_0(w)\subseteq\phi_0(z)
 \qquad(z,w\in S_0).
 \label{amp:reverse}
\end{equation}
Then
\begin{equation}
 2m\geq 2^{\lfloor(h-2)/127\rfloor}.
 \label{amp:bound}
\end{equation}
The map need not initially be assumed to preserve identities.
\end{theorem}

The proof is completed in Section~\ref{amp:section}. We first explain
the automata reduction, using the finite-computation representation
below. Its construction, proved in
Section~\ref{auto:finite-computations}, uses one additional state to
represent success by an exit past the right endmarker. Local diagram
edges record finitely many stay moves followed by one move out of a
cell; no bound on repeated crossings is imposed.

\begin{lemma}\label{auto:representation}
Let $B$ be an $s$-state 2NFA over a finite alphabet $\Sigma$ in the model
fixed in the introduction. There are a monoid homomorphism
\[
  \tau:\Sigma^*\longrightarrow\mathcal T_{s+1},
\]
two fixed endmarker diagrams $\lambda,\rho\in\mathcal T_{s+1}$,
and fixed positive labels $a,b$ such that, for every word $w$,
\begin{equation}\label{auto:recognition}
  w\in L(B)
  \quad\Longleftrightarrow\quad
  (a,b)\in F_{\lambda\tau(w)\rho}.
\end{equation}
Consequently, if $\tau(u)\sqsubseteq\tau(v)$, then
\begin{equation}\label{auto:context-monotonicity}
  \alpha u\beta\in L(B)
  \quad\Longrightarrow\quad
  \alpha v\beta\in L(B)
  \qquad(\alpha,\beta\in\Sigma^*).
\end{equation}
\end{lemma}
The acceptance equivalence immediately implies the stated monotonicity
in arbitrary word contexts:
\begin{proof}[Derivation of context monotonicity]
Diagram multiplication is increasing in each factor.
If $\tau(u)\sqsubseteq\tau(v)$, then
\[
 \lambda\tau(\alpha)\tau(u)\tau(\beta)\rho
 \ \sqsubseteq\
 \lambda\tau(\alpha)\tau(v)\tau(\beta)\rho.
\]
The designated forward edge persists under this inclusion.  Applying
\eqref{auto:recognition} proves
\eqref{auto:context-monotonicity}.
\end{proof}

For relation-product emptiness, singleton contexts convert this
increasing behavior of acceptance into a decreasing relation image.
For $F\subseteq H$, write $I_F=\{(x,x):x\in F\}$.
A multiplicative map is \emph{unital} if it preserves identities.

\begin{proposition}\label{auto:order-reversal}
Let $H$ be a finite set with at least two elements. If an $s$-state
2NFA recognizes $\Sigma_H^*\setminus L_H$, then a
submonoid of $\mathcal T_{s+1}$ admits a surjective unital
homomorphism $\phi$ onto $\mathcal R_H$ satisfying
\begin{equation}\label{auto:reverse-order}
  z\sqsubseteq z'
  \quad\Longrightarrow\quad
  \phi(z')\subseteq\phi(z).
\end{equation}
\end{proposition}
\begin{proof}
Apply Lemma~\ref{auto:representation} to the complementing machine,
and let $S=\tau(\Sigma_H^*)$.  Suppose
$\tau(u)\sqsubseteq\tau(v)$.  For each $p,q\in H$, take the
one-letter words
\[
  \alpha=I_{\{p\}},\qquad \beta=I_{\{q\}}.
\]
They belong to $\Sigma_H$, and
\[
 I_{\{p\}}r(u)I_{\{q\}}
 =
 \begin{cases}
   \{(p,q)\},& (p,q)\in r(u),\\
   \varnothing,&(p,q)\notin r(u).
 \end{cases}
\]
Since the machine accepts exactly the words with empty relation
product, \eqref{auto:context-monotonicity} implies
\[
  (p,q)\notin r(u)
  \quad\Longrightarrow\quad
  (p,q)\notin r(v).
\]
Thus
\begin{equation}\label{auto:word-order}
  \tau(u)\sqsubseteq\tau(v)
  \quad\Longrightarrow\quad
  r(v)\subseteq r(u).
\end{equation}

If $\tau(u)=\tau(v)$, applying \eqref{auto:word-order} in both
directions yields $r(u)=r(v)$.  Therefore
\[
  \phi:S\longrightarrow\mathcal R_H,
  \qquad \phi(\tau(w))=r(w),
\]
is well-defined.  Concatenation of words proves multiplicativity, and
the empty word proves preservation of the identity.  It is surjective
because every relation in $\mathcal R_H$ is itself a letter of
$\Sigma_H$.  Equation~\eqref{auto:word-order} is exactly the required
order reversal \eqref{auto:reverse-order}.
\end{proof}

With $m=s+1$, Theorem~\ref{amp:image} therefore bounds the states of a
complementing machine. Section~\ref{auto:section} constructs the
$h+2$-state source machine and completes the numerical deduction.
It remains to prove the algebraic bound. The next subsection associates
at most $2m$ recurrent classes with an idempotent diagram. We will measure
loss by the classes whose through loops are missing, and force that
loss to grow exponentially with $h$.

\subsection{Recurrent classes in a corner}

For an idempotent $e\in\mathcal T_m$, define
\begin{equation}\label{diag:KP}
\begin{aligned}
 K_e^+&=(L_eR_e)^*,& P_e^+&=K_e^+F_eK_e^+,\\
 K_e^-&=(R_eL_e)^*,& P_e^-&=K_e^-B_eK_e^-.
\end{aligned}
\end{equation}
When working only with the positive sign we omit superscripts.
By $e^2=e$ and \eqref{diag:product},
\begin{equation}\label{diag:idempotent}
 F_eK_eF_e=F_e,\qquad P_e^2=P_e.
\end{equation}
In particular, $P_e$ is transitive. A point $x\in D^\sigma$ is
\emph{recurrent} if $xP_e^\sigma x$. Mutual relatedness under
$P_e^\sigma$ is an equivalence relation on the recurrent points.
Let $\mathcal C_e$ be the collection of its classes for both signs,
keeping the signs distinct. Then
\begin{equation}\label{diag:class-count}
 |\mathcal C_e|\le2m.
\end{equation}
For the identity diagram, $K_e^\sigma=P_e^\sigma$ is the identity
relation on $D^\sigma$, so all $2m$ directional labels form singleton
recurrent classes.

The \emph{corner} $e\mathcal T_me=\{eae:a\in\mathcal T_m\}$ is a monoid
with identity $e$; membership is equivalent to $ez=ze=z$.
For $z$ in this corner put
\[
 U_e^+(z)=K_e^+F_zK_e^+,\qquad
 U_e^-(z)=K_e^-B_zK_e^-.
\]
The return formulas and the two corner identities give
\begin{equation}\label{diag:corner}
 L_e\subseteq L_z,\quad R_e\subseteq R_z,\qquad
 F_z=F_e(L_zR_e)^*F_z=F_z(L_eR_z)^*F_e.
\end{equation}
Consequently
\begin{equation}\label{diag:corner-absorb}
 F_eK_eF_z\subseteq F_z,\quad F_zK_eF_e\subseteq F_z,
 \qquad P_eU_e(z),\ U_e(z)P_e\subseteq U_e(z).
\end{equation}
The last two inclusions follow by multiplying the first two by $K_e$
on the appropriate sides. If $x,y$ belong to one recurrent class, then
$yP_ex$, $xP_ey$, and \eqref{diag:corner-absorb} show that a loop
$xU_e(z)x$ implies $yU_e(z)y$. Reflection gives the same fact for the
negative sign. Hence the following is well-defined:
\[
 \mathcal M_e(z)=\{C\in\mathcal C_e:
       xU_e^\sigma(z)x\text{ fails for }x\in C,
       \text{ where }C\subseteq D^\sigma\}.
\]
We call $\mathcal M_e(z)$ the \emph{missing set} of $z$ relative to $e$.

\begin{lemma}[Products of missing sets]\label{diag:missing}
Let $e\in\mathcal T_m$ be idempotent. For $z,w\in e\mathcal T_me$,
\[
 \mathcal M_e(zw)\subseteq\mathcal M_e(z)\cup\mathcal M_e(w),
 \qquad \mathcal M_e(e)=\varnothing.
\]
In particular, $\mathcal M_e(z^k)\subseteq\mathcal M_e(z)$ for every
positive integer $k$.
\end{lemma}
\begin{proof}
The returns of $e$ persist in both $z$ and $w$, so
$(L_wR_z)^*\supseteq K_e$. Formula~\eqref{diag:product} therefore gives
$U_e^+(zw)\supseteq U_e^+(z)U_e^+(w)$. Two loops at the same point
compose to a loop. For the negative sign the corresponding product is
$U_e^-(w)U_e^-(z)$, which gives the same conclusion. Finally
$U_e^\sigma(e)=P_e^\sigma$, so no recurrent class is missing from $e$.
\end{proof}

\begin{lemma}[Rectangles of through edges]\label{diag:rectangles}
Let $e\in\mathcal T_m$ be idempotent and let $z\in e\mathcal T_me$.
For a positive recurrent class $C\in\mathcal C_e$ and any $p\in C$, set
\[
 A_C=\{u:u(F_eK_e)p\},\qquad
 B_C=\{v:p(K_eF_e)v\}.
\]
The rectangle $A_C\times B_C$ is independent of $p$ and is contained
in $F_e$. These rectangles cover $F_e$. If
$C\notin\mathcal M_e(z)$, its rectangle is contained in $F_z$.
The reflected assertions hold for negative classes and $B_e$.
In particular,
\begin{equation}\label{diag:detect}
 \mathcal M_e(z)=\varnothing\quad\Longrightarrow\quad e\sqsubseteq z.
\end{equation}
\end{lemma}
\begin{proof}
Containment in $F_e$ follows from $F_eK_eF_e=F_e$.
The identities
\[
 (F_eK_e)P_e=F_eK_e,\qquad P_e(K_eF_e)=K_eF_e
\]
show that replacing $p$ by a mutually $P_e$-related point does not
change either factor of the rectangle.

To prove coverage, \eqref{diag:idempotent} implies
$F_eP_e^kF_e=F_e$ for every $k\ge1$. Given a pair in $F_e$, choose a
witness for this expression with $k>m$. Some label repeats along the
$P_e$ portion. That label $p$ is recurrent, since a nonempty $P_e$-path
from $p$ back to itself is a $P_e$-edge by transitivity. The portions
before and after $p$ belong to $F_eK_e$ and $K_eF_e$: indeed
$F_eP_e^j\subseteq F_eK_e$ and $P_e^jF_e\subseteq K_eF_e$ for
$j\ge0$, including $j=0$ since $K_e$ is reflexive.

If $C$ is not missing from $z$, insert $pK_eF_zK_ep$ between the two
halves of its rectangle. The resulting relation is contained in
$F_eK_eF_zK_eF_e\subseteq F_z$ by
\eqref{diag:corner-absorb}. Reflection proves the backward assertion.
When the missing set is empty all through edges of $e$ are thus present
in $z$, and its return edges persist by \eqref{diag:corner}.
\end{proof}

\subsection{Transport through a fixed context}

The next statement controls a whole family of replacements at once.
The common target set, rather than merely a bound for each replacement,
will be essential in the amplification argument.

\begin{lemma}[Transport]\label{diag:transport}
Let $e,d\in\mathcal T_m$ be idempotents and let $u,v\in\mathcal T_m$
satisfy $uev=d$. For every $J\subseteq\mathcal C_e$ there exists a single
set $J'\subseteq\mathcal C_d$, with $|J'|\le2|J|$, such that
\[
 \mathcal M_d(uzv)\subseteq J'
\]
for every $z\in e\mathcal T_me$ satisfying
$\mathcal M_e(z)\subseteq J$ and $uzv\in d\mathcal T_md$.
\end{lemma}
\begin{proof}
For each positive class of $\mathcal C_d$, choose a representative $x$
and a witness for $xK_d^+F_{uev}K_d^+x$. Expand the middle edge into an
actual finite path through the three diagrams $u,e,v$. For a negative
class do the same with $xK_d^-B_{uev}K_d^-x$. Make these choices once,
independently of $z$.

For every occurrence of a through edge of the middle factor $e$ on a
chosen path, choose an $e$-class whose rectangle covers that edge, using
Lemma~\ref{diag:rectangles}. Mark a $d$-class if any class chosen along
its witness lies in $J$; let $J'$ be the marked set. If a class is
unmarked, all its chosen through edges survive replacement of $e$ by
$z$, as do all return edges of $e$. Its original outer $K_d^\sigma$
pieces remain fixed. Thus its witness is a loop in $U_d^\sigma(uzv)$,
and the class is not missing. This proves the required common-set
inclusion.

It remains to count marked classes. For each $d$-class $C$, let $S_C$
be the set of $e$-classes chosen along its fixed witness. We show that
the sets $S_C$ are pairwise disjoint among $d$-classes of a fixed sign.
Consider two positive witnesses, written as
\[
 xK_d^+\alpha\;F_d\;\beta K_d^+x,
 \qquad yK_d^+\alpha'\;F_d\;\beta'K_d^+y,
\]
where $\alpha,\beta,\alpha',\beta'\in D^+$ and the displayed
$F_d$-edges have been expanded through $u,e,v$.
Suppose the expanded paths use occurrences of middle-factor edges
$(a,b)$ and $(a',b')$ covered by one $e$-class. Its rectangle also
contains $(a,b')$ and $(a',b)$. The sign of that $e$-class fixes the
orientation of both edges and their entry and exit boundaries inside
the three-factor product. Thus the prefix ending at $a$, the crossed
edge $(a,b')$, and the suffix starting at $b'$ concatenate to a path
from $\alpha$ to $\beta'$. Neither retained subpath has an earlier
outer exit. The concatenation is therefore a valid finite path for
$F_d$, even if it repeats internal crossings or its selected middle
edge points backward. Interchanging the two prefixes and suffixes
similarly gives $\alpha'F_d\beta$.

The original exterior factors now give $xP_d^+y$ and $yP_d^+x$, so
$x$ and $y$ belong to the same recurrent class. Reflection gives the
same conclusion for negative witnesses, with $B_d$ in place of $F_d$.
This proves the asserted disjointness for each sign. A witness using
only returns of $e$ has $S_C=\varnothing$ and is never marked. Since
$J'=\{C:S_C\cap J\ne\varnothing\}$, at most $|J|$ classes of each
sign are marked, and $|J'|\le2|J|$.
\end{proof}
